% -------------------------------------------------------------------------
% WBS ID: RES-VER-FRM
% Deliverable: Simulation Credibility and Test Framework
% Status: In progress
% Evidence status: Local 3D validation hierarchy established
% Review status: Not reviewed
% Last updated: 2026-07-19
% Dependencies: RES-MOD-SPEC, RES-NUM-SPEC, RES-PHY-LOD
% Downstream interfaces: RES-VER-BMK, RES-VER-MET, Software validation reports
% -------------------------------------------------------------------------

\section{RES-VER-FRM --- Simulation Credibility and Test Framework}
\label{sec:res-ver-frm}

\subsection{Purpose and Scope}
\label{sec:res-ver-frm-purpose}

% TODO: Define what this Deliverable covers and explicitly excludes.

This Deliverable defines the validation hierarchy for the local
three-dimensional URANS--VOF tsunami--barrier model. It separates
solver capability, event-conditioned hydrodynamics and local
impact/loading credibility.

\subsection{Research Questions}
\label{sec:res-ver-frm-research-questions}

% TODO: State the questions that must be answered before the Deliverable
% can be accepted.

\subsection{Terminology and Definitions}
\label{sec:res-ver-frm-definitions}

% TODO: Define the technical terminology, symbols, quantities and
% classifications used in this Deliverable.

\subsection{Literature Evidence}
\label{sec:res-ver-frm-literature-evidence}

% TODO: Record evidence from peer-reviewed papers, authoritative datasets,
% standards, technical reports and primary documentation.
%
% Do not create a detached literature-summary list. Explain how each source
% supports, contradicts or limits a project decision.

\textcite{qin2018comparison} is retained as evidence that local model
validation must include velocity and force, not only free-surface
height. \textcite{watanabe2022validation} is retained as evidence that
blind tsunami-impact simulations can vary materially in force and
velocity predictions even when the general hydrodynamic setup is
shared. Source role: `3D-VALIDATION`. Project conclusion: local
validation shall be staged and multi-output. Limitation: passing one
stage does not automatically validate the next stage.

\subsection{Governing Relations and Quantitative Definitions}
\label{sec:res-ver-frm-governing-relations}

% TODO: Record relevant equations, dimensionless groups, empirical models,
% extraction rules or quantitative definitions.
%
% Use align environments for displayed equations.
% Every displayed equation must have a unique \label{eq:...}.
% Do not place punctuation after displayed equations.

\subsection{Test Objective}
\label{sec:res-ver-frm-test-objective}

% TODO: Complete this RES-VER Domain-specific subsection.

The local validation objective is:

\begin{quote}
  V1 solver capability
  $\rightarrow$
  V2 event-conditioned hydrodynamic realism
  $\rightarrow$
  V3 local impact and loading credibility
\end{quote}

V1 checks that the numerical method can solve canonical free-surface
and impact problems. V2 checks that the selected event and local
forcing reproduce relevant hydrodynamic observations. V3 checks that
pressure, force, impulse, moment, overtopping and transmitted/reflected
wave outputs are credible for barrier-performance assessment.

\subsection{Reference or Benchmark Solution}
\label{sec:res-ver-frm-reference-benchmark-solution}

% TODO: Complete this RES-VER Domain-specific subsection.

\subsection{Test Configuration}
\label{sec:res-ver-frm-test-configuration}

% TODO: Complete this RES-VER Domain-specific subsection.

\subsection{Measured Quantities}
\label{sec:res-ver-frm-measured-quantities}

% TODO: Complete this RES-VER Domain-specific subsection.

The framework shall measure free-surface elevation, arrival time,
velocity, run-up, overtopping, phase volume, pressure, shear, resultant
force, force impulse, moment, transmission/reflection measures and
computational diagnostics.

\subsection{Error Metrics}
\label{sec:res-ver-frm-error-metrics}

% TODO: Complete this RES-VER Domain-specific subsection.

\subsection{Tolerance and Acceptance Criteria}
\label{sec:res-ver-frm-tolerance-acceptance-criteria}

% TODO: Complete this RES-VER Domain-specific subsection.

Acceptance shall be stage-specific. A model that passes V1 remains
unvalidated for the selected event until V2 passes. A model that passes
V2 remains unvalidated for barrier loading until V3 passes. URANS
$k$--$\omega$ SST shall be accepted for screening only where it
preserves the safety-critical quantities and candidate ranking within
the tolerances defined in `RES-VER-MET`.

\subsection{Execution Handoff}
\label{sec:res-ver-frm-execution-handoff}

Research deliverables specify the credibility hierarchy and required
metrics only. Software validation reports shall state whether a result
is V1, V2 or V3 evidence, which dataset and geometry identifiers were
used, and which acceptance criteria passed. No validation shall be
claimed from the existence of this specification alone.

\subsection{Candidate Approaches}
\label{sec:res-ver-frm-candidate-approaches}

% TODO: Define the credible candidate approaches considered.

\subsection{Critical Comparison}
\label{sec:res-ver-frm-critical-comparison}

% TODO: Compare candidates using physically and project-relevant criteria.
% Do not select an approach solely on popularity or implementation ease.

\subsection{Assumptions and Applicability}
\label{sec:res-ver-frm-assumptions}

% TODO: State assumptions and the conditions under which they are valid.

\subsection{Limitations and Failure Conditions}
\label{sec:res-ver-frm-limitations}

% TODO: State where the evidence, model or selected approach ceases to be
% reliable or applicable.

\subsection{Project-Specific Conclusions}
\label{sec:res-ver-frm-conclusions}

% TODO: Convert the evidence into conclusions specific to the Studentship
% project. Do not merely restate literature findings.

\subsection{Selected Approach and Rationale}
\label{sec:res-ver-frm-selected-approach}

% TODO: Record the selected approach only after sufficient evidence exists.
% State the alternatives rejected and the reasons for rejection.

The selected credibility framework is the three-stage V1--V2--V3
hierarchy. It prevents water-level agreement, solver stability or
successful regional propagation from being treated as sufficient
validation of local pressure and loading.

\subsection{Unresolved Decisions}
\label{sec:res-ver-frm-unresolved-decisions}

% TODO: Record unresolved questions, required evidence and decision owners.

\subsection{Requirements and Handoffs}
\label{sec:res-ver-frm-requirements}

% TODO: State requirements passed to other Research Domains, Software,
% verification activities or later execution work.

`RES-VER-BMK` shall define benchmark classes for V1. `RES-VER-MET`
shall define output-specific metrics and tolerances for V2 and V3.
Software validation reports shall identify which stage each result
supports.

\subsection{Deliverable Completion Record}
\label{sec:res-ver-frm-completion-record}

% TODO: Record whether the Deliverable is:
%
% - Not started
% - In progress
% - Draft complete
% - Reviewed
% - Accepted
%
% Acceptance requires:
%
% - the research questions to be answered;
% - claims to be supported by appropriate evidence;
% - assumptions and limitations to be explicit;
% - the project-specific conclusion to be stated;
% - downstream requirements to be traceable;
% - unresolved decisions to be closed or formally deferred.
